\section*{4. Experiments}
\subsection*{4.1 Tasks, datasets and evaluation}

We evaluate BioCoLoop on three biological prediction problems. Each task uses its own base model, while collaborative training, evaluation, and proposal generation follow the same procedure. We consider two ways of forming laboratories. In the main setting, each dataset is divided into ten disjoint groups. In a complementary setting, distinct biological datasets or experimental scenarios are treated as separate laboratories. The first setting provides a controlled comparison of collaborative training, while the second tests whether information can transfer across genuinely different data sources.

\textbf{Drug-target interactions.}
Given a compound and a protein sequence, the model predicts whether they interact. BindingDB provides 50,149 training pairs and 5,604 development pairs, with canonical compounds defining laboratory assignments. We use TAPB as the base model, with frozen ESM2 protein features and trainable drug, fusion, and classification modules \citep{lin2025tapb}. The primary metric is AUROC on 5,505 random test pairs. Additional test sets containing 3,791 unseen-drug pairs and 2,507 unseen-protein pairs measure generalization to unseen entities. AP, MCC, accuracy, and cross-entropy are reported as secondary metrics.

\textbf{Observed-response proteomics.}
Given a drug representation and its measured 6h and 24h proteomic responses, the model predicts binary efficacy. A ProteinTalks PTPC subset provides 386 training, 105 development, and 148 test conditions \citep{Sun2026ProteinTalks}. Compound-group assignment gives each laboratory 49-50 training and development conditions. The base model is a ProteinTalks-derived efficacy model with protein and drug branches, trained from random initialization. Average precision (AP) is the primary metric. AUROC, accuracy, and cross-entropy are secondary metrics. Unlike de novo response prediction, this task conditions on observed post-treatment proteomic measurements.

\textbf{Cell perturbation identification.}
Given an expression response and five candidate interventions, the model identifies the intervention whose predicted response has the highest cosine similarity to the observation. We use VCC single-gene perturbations \citep{arc2025vcc}, Norman double-gene perturbations \citep{norman2019genetic}, and a Tahoe subset of drug perturbations in A549 at 5 $\mu$M \citep{zhang2025tahoe}. Intervention-disjoint training, development, and test counts are 60/20/20, 70/20/27, and 180/60/67, respectively. The base model uses a frozen scDEBART backbone with a trainable response head \citep{sung2026scdebart}. All methods use the same corrected attention masks and the same candidate sets. The primary metric is intervention-macro Top-1, which first averages accuracy within each intervention and then across interventions. Random selection has an expected Top-1 of 20\%. Micro Top-1, MRR, Top-3, and cross-entropy are reported as secondary metrics. Appendix A provides the response construction, chemical conditioning, training losses, and secondary evaluation details.

\subsection*{4.2 Compared methods and shared protocol}

\textbf{Compared methods.}
The \emph{fixed model} trains the task-specific base model with a fixed training recipe. \emph{Direct optimization} uses a self-improving model to propose changes to the model structure or training hyperparameters without access to results from previous experiments. AI-Scientist-v2 and AI-Researcher are adapted to the same biological tasks and evaluate their proposed changes through the same training and scoring procedure. BioCoLoop additionally provides results from completed experiments when generating subsequent proposals and performs collaborative training and evaluation across laboratories. We repeat the comparison with Qwen2.5-7B-Instruct and GPT-5.6 Luna as proposal generators. Details of the AI-Scientist-v2 and AI-Researcher adaptations are given in Appendix C.

The fixed, direct, AI-Scientist-v2, and AI-Researcher baselines use laboratory 0 because these methods do not support collaborative training and evaluation. BioCoLoop uses ten laboratories. The main table therefore compares complete systems rather than isolating the effect of laboratory access. To separate these effects, Appendix A reports a matched comparison that crosses fixed, direct, and BioCoLoop-style proposal generation with one or ten laboratories.

\textbf{Shared protocol.}
Direct optimization and BioCoLoop use the same twelve predefined combinations of model structure and training hyperparameters, the same six proposals, and the same first proposal. Each proposed model starts from the same task and seed-specific initialization and is trained for 100 rounds. Evaluation is performed every five rounds. Model selection averages the primary metric equally across participating laboratories. We report mean and sample standard deviation over seeds 42-44 using fixed data partitions. All methods within a task use the same held-out test set and evaluation code. Appendix A provides the complete training settings, matched comparisons, and uncertainty estimates.

% \subsection*{4.1 Tasks, datasets and evaluation}
% We study three biological prediction problems, using a separate predictor for each task and the same BioCoLoop research interface. The main comparison uses ten simulated laboratories formed by disjoint groups within each dataset. A complementary source-as-laboratory study assigns different biological collections or scenarios to separate laboratories. The first setting controls access to data; the second tests transfer between data sources.

% \paragraph{Drug--target interactions.}
% Given a compound and a protein sequence, the predictor estimates whether they interact. BindingDB supplies 50,149 training and 5,604 development pairs; canonical compounds define the laboratory assignments. The reference is TAPB with frozen ESM2 protein features and trainable drug, fusion and classification modules \citep{lin2025tapb}. The main endpoint is AUROC on 5,505 random-test pairs. Tests with 3,791 unseen-drug and 2,507 unseen-protein pairs assess entity generalization. AP, MCC, accuracy and cross-entropy provide secondary measures.

% \paragraph{Observed-response proteomics.}
% Given a drug representation and its measured 6h/24h proteomic response, the predictor estimates binary efficacy. A ProteinTalks PTPC subset supplies 386 training, 105 development and 148 test conditions \citep{Sun2026ProteinTalks}. Compound-group assignment gives each laboratory 49--50 training-plus-development conditions. The reference is a ProteinTalks-derived efficacy head with protein and drug branches, fitted from random initialization. Average precision (AP) is primary; AUROC, accuracy and cross-entropy are secondary. This task uses observed post-treatment measurements, rather than predicting efficacy from an unmeasured proteome.

% \paragraph{Cell perturbation identification.}
% Given an expression response and five candidate interventions, the model identifies the intervention whose predicted response has the highest cosine agreement with the observation. We use VCC single-gene perturbations \citep{arc2025vcc}, Norman double-gene perturbations \citep{norman2019genetic}, and a Tahoe subset of drug perturbations in A549 at 5 $\mu$M \citep{zhang2025tahoe}. Intervention-disjoint training/development/test counts are 60/20/20, 70/20/27 and 180/60/67, respectively. The reference is a frozen scDEBART backbone with a trainable response head \citep{sung2026scdebart}. All arms use the same corrected attention-mask implementation and fixed candidate rosters. Primary scoring is intervention-macro Top-1: accuracy is first averaged within each intervention, then across interventions. Random choice has expectation 20\%. Micro Top-1, MRR, Top-3 and cross-entropy are secondary. Appendix A describes response construction, chemical conditioning, losses and secondary metrics.

% \subsection*{4.2 Compared methods and shared protocol}
% The \emph{fixed model} trains each reference architecture with a fixed recipe on preselected laboratory 0. \emph{Direct optimization} proposes recipes without evaluation history, while AI-Scientist-v2 and AI-Researcher are task-adapted re-implementations that keep their native proposal logic but route candidates through the registered design interface on laboratory 0, with code execution, retrieval and manuscript generation disabled (Appendix~C). \emph{BioCoLoop} supplies evidence-card feedback to the proposer while training across ten laboratories. The main comparison repeats these methods in two research-model blocks, local Qwen2.5-7B-Instruct and GPT-5.6 Luna at low reasoning effort. A six-arm factorial comparison crosses fixed/direct/loop research with one/ten laboratories to separate data access from proposal feedback.

% Direct and loop search share a twelve-design library, six proposal slots and the first proposal. Every candidate starts from the same task/seed initialization and receives 100 training rounds, with development evaluation every five rounds. Selection uses equal-weight development scores over participating laboratories. The main results report mean and sample SD over seeds 42--44 on fixed data partitions. All methods within a task share the held-out test set and scorer. Appendix A gives the full factorial comparison, complete optimization settings and uncertainty estimates.

% Task-adapted AI-Scientist-v2 and AI-Researcher \citep{Yamada2025AIScientistV2,Tang2025AIResearcher} access laboratory 0 because they do not implement collaborative fitting or distributed evaluation, while BioCoLoop uses ten. All rows share the design library, candidate budget, training schedule and scorer. Qwen decoding uses registered generation seeds; the Luna service exposes no generation seed, so its seed labels identify training/search repetitions rather than deterministic sampling. Appendix~A separates loop and participation effects; Appendix~C details the adaptations.
